\chapter{Analiz Modeli}

\section{Genel Bakış}
Analiz modeli, sistemin iş mantığını ve kullanıcı etkileşimlerini anlamak için oluşturulmuştur. UML diyagramları kullanılarak gereksinimler görselleştirilmiştir.

\section{Use Case Diyagramı}
Sistemin aktörleri ve fonksiyonel gereksinimlerin etkileşimi Şekil \ref{fig:usecase}'de gösterilmiştir.

\begin{figure}[H]
    \centering
    \includegraphics[width=0.9\textwidth]{use-case.png}
    \caption{Use Case Diyagramı}
    \label{fig:usecase}
\end{figure}

\textbf{Aktörler:} Kimlik Doğrulanmış Kullanıcı (Authenticated User).

\textbf{Temel Senaryolar:} Kayıt olma, giriş yapma, mesajlaşma, dosya paylaşımı ve grup yönetimi gibi temel işlevleri kapsar.

\section{Domain Model (Etki Alanı Modeli)}
Sistemin temel varlıkları ve bu varlıklar arasındaki ilişkiler Şekil \ref{fig:domain}'de sunulmuştur.

\begin{figure}[H]
    \centering
    \includegraphics[width=0.9\textwidth]{domain-model.png}
    \caption{Domain Model Diyagramı}
    \label{fig:domain}
\end{figure}

\textbf{Ana Varlıklar:}
\begin{itemize}
    \item \textbf{User:} Kullanıcı kimlik ve profil bilgileri.
    \item \textbf{Chat:} Sohbet odası (birebir veya grup).
    \item \textbf{Message:} İçerik, tür ve durum bilgisi.
    \item \textbf{File:} Dosya metadata bilgileri.
    \item \textbf{Notification:} Bildirim detayları.
\end{itemize}

\section{Activity Diyagramları}
İş süreçlerinin akışını gösteren aktivite diyagramları aşağıda yer almaktadır.

\begin{figure}[H]
    \centering
    \includegraphics[width=0.8\textwidth]{activity-diagram.png}
    \caption{Mesaj Gönderme Aktivite Diyagramı}
    \label{fig:activity-msg}
\end{figure}

\begin{figure}[H]
    \centering
    \includegraphics[width=0.8\textwidth]{activity-diagram-file-upload.png}
    \caption{Dosya Yükleme Aktivite Diyagramı}
    \label{fig:activity-file}
\end{figure}
